\documentclass[10pt,a4paper]{amsart}
\usepackage[margin=19mm]{geometry}
\usepackage{amsmath,amssymb,mathtools}
\usepackage[scr=rsfs,cal=euler]{mathalfa}
\usepackage{xfrac}
\usepackage[unicode,hidelinks,bookmarks=false]{hyperref}
\newcommand{\ie}{i.e.\ }
\newcommand{\cf}{cf.\ }
\newcommand{\resp}{respectively\ }
\newcommand{\Subsection}{\subsection}
\providecommand{\nobreakdash}{\nobreak\hbox{-}}
\setlength{\emergencystretch}{2em}
\multlinegap0pt
\usepackage{bm}
\usepackage{bbm}
\usepackage{dsfont}
\usepackage{xspace}
\usepackage{cancel}
\usepackage{mathtools}
\usepackage{amsthm}
\makeatletter
\@ifundefined{lemma}{\newtheorem{lemma}{Lemma}}{}
\makeatother
\newcommand{\cC}{\mathcal{C}}
\title[Slow Convergence of Ising and Spin Glass Models with Well-Se]{Slow Convergence of Ising and Spin Glass Models with Well-Separated Frustrated Vertices}
\author{David Gillman, Dana Randall}
\date{Source: arXiv:2605.25323v1 (CC BY 4.0)}
\begin{document}
\maketitle

\noindent\emph{Source and licence.} Slow Convergence of Ising and Spin Glass Models with Well-Separated Frustrated Vertices. Version arXiv:2605.25323v1 (CC BY 4.0), distributed under the Creative Commons Attribution 4.0 International licence, as recorded in the arXiv licence metadata for this exact version and shown on the abstract page https://arxiv.org/abs/2605.25323v1. Full rights notice: licenses/arxiv-2605.25323-CC-BY-4.0.txt.\par\smallskip

\noindent\textit{Editorial scope.} Counted unit: the Lemma whose source label is \texttt{lem:piOmegaL1} in the article (source0.tex lines 541--547, source label \texttt{lem:piOmegaL1}), delivered together with the article's own complete original proof of it (source0.tex lines 549--579, one \texttt{proof} environment of 31 lines). No mathematics is written, repaired or re-derived: statement and proof are the article's own text line for line. Required context retained uncounted: lem:piOmegaC (lines 505--512). Every definition, display and label of those lines is retained unchanged. Omitted from this package and not redistributed: 1--504, 513--540, 548--548, 580--999. Nothing inside a retained excerpt is omitted or altered: every retained body is delivered exactly as the article's lines are written, so no omission receipt of any kind accompanies this package. Citations: the retained excerpts carry no citation command at all, so no bibliography entry of the article is transcribed and no reference is redistributed. Reference adaptation: none was needed and none was made; every reference inside the delivered unit resolves inside it, so no unresolved reference or citation is produced. Printed slips kept literal: no printed word, symbol, hypothesis or number of the counted statement or of its proof was altered. Layout and class: the article's own class is \texttt{amsart} with packages including microtype, amsmath, amsfonts, amssymb, xcolor, epsfig, fontenc, verbatim, gensymb, mathrsfs, amsbsy, comment, textcomp, framed; the wrapper is the lane's pinned \texttt{amsart} editorial wrapper at 10pt on A4, which restates the theorem-like environments the retained text needs and adds the article's own macro definitions that the retained text needs (\texttt{\textbackslash cC}), with extra wrapper packages (bm, bbm, dsfont, xspace, cancel, mathtools). The delivered numbering is the unit's own. Assets: no retained span contains a figure environment, a graphics-inclusion command, a PDF-page inclusion, a table or a TikZ picture; no image file is copied into this package, the package declares no asset dependency and no image file is redistributed with it. Licence: version arXiv:2605.25323v1 of this article is distributed under the Creative Commons Attribution 4.0 International licence, as recorded in the arXiv licence metadata for this exact version and shown on the abstract page https://arxiv.org/abs/2605.25323v1. A full rights notice is delivered at licenses/arxiv-2605.25323-CC-BY-4.0.txt. Characters: every retained excerpt is pure ASCII, so the delivered engine, XeLaTeX with its default Latin Modern text font, reports no missing character.
\begin{lemma}\label{lem:piOmegaC}
For any cross $C$ of length $|C| = L + \ell$   we have
\begin{equation*}
\pi(\cC_C ) \leq  \lambda^{-(4n-8np+\ell/2)} 2^{-f(\ell, p)},
\end{equation*}
where 
$\phi = {(1+\sqrt{5})}/{2}$.
\end{lemma}

\begin{lemma}\label{lem:piOmegaL1}
Let $\cC_{\ell}$ be the set of contour configurations containing a cross of length $L + \ell$. 
Then for large enough $\lambda$ 
\begin{equation}\label{eqn:piOmega0L}
- \log_2(\pi(\cC_{\ell}))/n = \Omega(1).
\end{equation}
\end{lemma}

\begin{proof}
Let $s = 1/2 - p$, $r = \ell/n$. Abusing notation, let $f(r, s) = f(\ell, p)/n$.

As $s, r \rightarrow 0$, $f(r, s) \rightarrow 4 - 4\log_2\phi \approx 1.2$.

Each $C \in \cC_{\ell}$ contains a minimal cross consisting of a vertical path connecting $v_1$ to $v_3$ and a horizontal path connecting $v_2$ to $v_4$. 
The vertical path contains $2n$ vertical edges and $2n-4pn$ horizontal edges. 
There are at most $\binom{4n-4pn}{2n-4pn} = \binom{2n+4sn}{4sn}$ choices of  vertical path.  
Likewise there are at most $\binom{2n+4sn}{4sn}$ choices of minimal horizontal path.  
Then there are $\binom{8n-8np+\ell}{\ell} = \binom{4n+8sn+rn}{rn}$ ways to choose the locations of the $\ell$ extra edges, and $3$ possible directions for each extra edge. In sum,

\begin{align*}
\pi(\cC_{\ell}) & \leq | \{C : |C| = L + \ell \}| \max_{|C| = L + \ell} \pi(\cC_C ) \\
      & \leq \binom{4n+8sn+rn}{rn} \binom{2n+4sn}{4sn}^2 3^{rn} 
      \max_{|C| = L + \ell} \pi(\cC_C ) \\
      & \leq \exp_2\left[n(4+8s+r)h\left(\frac{r}{4+8s+r}\right) + n(4+8s)h\left(\frac{2s}{1 + 2s}\right)\right] 3^{rn} \\
      &  \qquad \cdot \max_{|C| = L + \ell} \pi(\cC_C ) \\
      & \leq \exp_2[n(~ (4+8s+r)h(r/(4+8s+r)) + (4+8s)h(2s/(1 + 2s)) + r\log_2 3  \\
      & \qquad\qquad  - \log_2\lambda(8s + r/2) - f(r, s)~)] ,
\end{align*}
where we have applied first Stirling's formula and then Lemma~\ref{lem:piOmegaC} in the last two inequalities.

Taking $\log_2$ of both sides, dividing by $n$, and rearranging terms, we have
\begin{align*}
- \log_2(\pi(\cC_{\ell}))/n & \geq  
                -[~(4+8s+r)h(\frac{r}{4+8s+r}) + (4+8s)h(\frac{2s}{1+2s}) + r \log_2 3 ~] \cr
				     & \qquad + [~ (8s+ r/2) \log_2 \lambda + f(r, s) ~] .
\end{align*}

For small $s$ and $r$ the first term is $o(1)$ and the second term is $\Omega(1)$. Otherwise, for large $\lambda$ the second term dominates.
\end{proof}

\end{document}
